% Vietnamese translation for OI Public Library
% Source-PDF: IOI中国国家候选队论文/2001/高岳.pdf
\documentclass[11pt]{article}
\usepackage{oipl-vn}

\title{Báo cáo lời giải bài ``Độ cứng trung bình''}
\author{Cao Nhạc}
\date{}

\begin{document}
\maketitle

\origtitle{Medium Hardness solution report}
\sourcepath{IOI China candidate team papers / 2001 / Gao Yue.pdf}

\section*{Tóm tắt}

``Độ cứng trung bình'' là bài cuối của ngày thi thứ nhất IOI 2000. Bài này
chủ yếu kiểm tra khả năng sáng tạo của thí sinh: tự thiết kế một thuật toán
đúng và hiệu quả để giải quyết vấn đề. Bản báo cáo này trình bày quá trình và
phương pháp mà tác giả đã dùng khi làm bài.

\paragraph{Từ khóa.}
Tìm kiếm nhị phân; số ngẫu nhiên.

\section{Mô tả bài toán}

Xem phụ lục của đề gốc.

\section{Phân tích bài toán}

\subsection{Thuật toán 1-1}

Vì mỗi lần so sánh có thể xác định được vật lớn nhất và nhỏ nhất trong ba vật
(dù không biết phía nào là ``lớn'' về độ cứng tuyệt đối), ta có thể trước hết
tìm ra vật lớn nhất và nhỏ nhất trong các vật số $1,2,3$, rồi dùng hai vật đó
so với vật số $4$ để thu được vật lớn nhất và nhỏ nhất trong $1,\ldots,4$.
Tiếp tục như vậy với vật số $5$, rồi các vật còn lại, ta có thể tìm được hai
đầu mút của $1,\ldots,n$.

Hai vật ở hai đầu chắc chắn không phải vật có độ cứng trung bình, nên loại
chúng đi. Sau đó lại dùng phương pháp trên để tìm hai đầu mút trong $n-2$ vật
còn lại, rồi tiếp tục loại bỏ. Cuối cùng vật duy nhất còn lại chính là vật có
độ cứng trung bình.

Độ phức tạp số lần so sánh của thuật toán này là $O(n^2)$, không thỏa giới
hạn tìm trong $1449$ vật bằng $7777$ lần so sánh. Nguyên nhân là lượng thông
tin khai thác từ mỗi lần so sánh còn quá ít. Một lần so sánh ba vật thật ra
cho ta quan hệ thứ tự của ba vật đó, dù không biết chiều tăng hay giảm. Nếu đã
biết quan hệ thứ tự của hai vật, khi so sánh thêm với vật thứ ba ta biết được
vị trí của vật thứ ba trong thứ tự ấy. Thuật toán 1-1 không dùng thông tin này,
nên độ phức tạp cao. Từ góc nhìn sắp thứ tự đó, ta có thuật toán 2-1.

\subsection{Thuật toán 2-1}

Giả sử đã biết quan hệ thứ tự độ cứng của $m$ vật đầu tiên. Với vật tiếp theo,
dùng tìm kiếm nhị phân để chèn nó vào vị trí thích hợp, từ đó thu được quan hệ
thứ tự của $m+1$ vật đầu tiên. Lặp lại quá trình này cho đến khi có thứ tự của
tất cả $n$ vật, rồi lấy vật ở vị trí giữa.

Xét ví dụ sau:
\[
\begin{array}{c|ccccc}
\text{Nhãn} & 1 & 2 & 3 & 4 & 5\\
\hline
\text{Độ cứng} & 2 & 5 & 4 & 3 & 1
\end{array}
\]

\[
\begin{array}{c|c|c|c|c}
\text{Vật chèn} & \text{Thứ tự đã biết} & \text{So sánh}
& \text{Kết quả trả về} & \text{Thứ tự thu được}\\
\hline
& 123 & 3 & & 132\\
4 & 132 & 134 & 4 & 1( )32\\
5 & 1432 & 435 & 4 & (1)432\\
5 & 1432 & 145 & 1 & ()1432\\
& & & & 51432
\end{array}
\]

Dấu ngoặc biểu thị vùng mà vật mới có thể được chèn vào. Hàng thứ hai cho biết
vật $4$ nằm giữa $1$ và $3$. Hàng thứ ba cho biết vị trí của vật $5$ nằm ở bên
trái vật $4$, nên cần thêm một lần so sánh để xác định quan hệ giữa $5$ và $1$.

Thuật toán này dùng chèn nhị phân, với chi phí chèn một vật là
$O(\log_2 m)$. Cần chèn $n-3$ vật, nên tổng độ phức tạp nhỏ hơn phương pháp
thô. Thực nghiệm cho thấy khi $n=1449$ thì trung bình cần hơn $11000$ lần so
sánh, phù hợp với ước lượng
\[
  1449\log_2 1449 \approx 15216.
\]
Nguyên nhân là thuật toán đã sắp thứ tự tất cả vật, trong khi có nhiều vật
chắc chắn không thể là vật có độ cứng trung bình; việc sắp thứ tự chi tiết cho
chúng làm lãng phí số lần so sánh. Trên cơ sở thuật toán 2-1, ta thêm cắt tỉa
để được thuật toán 2-2.

\subsection{Thuật toán 2-2}

Đặt $2m+1=n$. Với một dãy đã được sắp thứ tự hoàn toàn, vật có độ cứng trung
bình có $m$ vật ở bên trái và $m$ vật ở bên phải. Vì thế, nếu biết một vật có
hơn $m$ vật nằm ở bên trái hoặc bên phải trong dãy hiện tại, vật đó chắc chắn
không phải đáp án. Khi ấy không cần xác định chính xác vị trí của nó; chỉ cần
chèn nó vào đầu trái hoặc đầu phải tương ứng.

Theo nguyên lý này, trước hết sắp thứ tự $m+1$ vật đầu tiên, vì mỗi vật trong
số đó đều còn có khả năng là đáp án. Khi vật thứ $m+2$ được chèn vào, hai vật
ở hai đầu chắc chắn không phải đáp án. Tương tự, khi hiện đã có $m+3$ vật được
sắp thứ tự, hai vật ở mỗi đầu của dãy cũng chắc chắn không phải đáp án. Nói
cách khác, vật thứ $k$ chỉ có thể trở thành vật trung bình nếu nó được chèn
vào đoạn vị trí từ $k-m-1$ đến $m+1$. Trong quá trình chèn nhị phân, nếu vùng
chèn hiện tại đã nằm ngoài đoạn có thể đó, vị trí cụ thể của vật mới không ảnh
hưởng đến kết quả, nên ta chèn thẳng nó vào đầu trái hoặc đầu phải.

Ví dụ đơn giản: giả sử đã sắp thứ tự $n-1$ vật. Trong số này chỉ hai vật ở
giữa còn có thể là đáp án; gọi chúng là $A$ và $B$, với $A$ ở bên trái $B$.
Theo cách trên, vật thứ $n$, gọi là $C$, chỉ cần so sánh một lần với $A$ và
$B$. Nếu $A$ là vật ở giữa trong phép so sánh, vùng chèn của $C$ nằm bên trái
$A$ và không còn nằm trong vùng có thể, nên chèn $C$ vào đầu trái; kết quả
cuối cùng là $A$. Nếu $B$ ở giữa, tương tự chèn $C$ vào đầu phải và kết quả là
$B$. Nếu $C$ ở giữa, chèn $C$ vào giữa $A$ và $B$ và kết quả là $C$. Không cần
như thuật toán 2-1, vốn phải so sánh thêm nhiều lần để tìm đúng vị trí của
$C$.

Cắt tỉa thứ hai dựa trên cắt tỉa thứ nhất. Thực nghiệm cho thấy khá nhiều vật
được chèn vào hai đầu nhờ cắt tỉa, nhưng trước khi biết chúng không thuộc vùng
có thể, thuật toán đã thực hiện $2$ đến $5$ lần phán đoán. Vì vậy có thể dùng
trước $1$ đến $2$ lần so sánh để kiểm tra vật cần chèn có nằm trong vùng có
thể hay không. Nếu có, mới dùng chèn nhị phân; nếu không, chèn trực tiếp vào
một trong hai đầu. Cách này cải thiện thêm một phần hiệu suất.

Cải tiến này đưa số lần so sánh trung bình cho $1449$ vật xuống khoảng
$8200$, chỉ còn cách giới hạn $7777$ một bước. Tuy vậy biên độ cải tiến rất
nhỏ, vì phương pháp sắp thứ tự này không cắt tỉa được việc sắp thứ tự $725$
vật đầu tiên, vốn đã tốn khoảng $5000$ lần so sánh. Do đó để giải bài trong
số lần quy định, cần thay đổi thuật toán và chuyển sang nhóm thuật toán 3-X.

\subsection{Thuật toán 3-1}

Tuy hai nhóm thuật toán trước không thành công, tác giả rút ra hai kinh
nghiệm đáng dùng tiếp:
\begin{enumerate}
  \item Cần có khái niệm vị trí. Dù không biết chiều lớn nhỏ tuyệt đối, vẫn
  phải giữ được thứ tự tương đối.
  \item Cần có khái niệm vùng. Trước hết xác định vùng có thể, rồi mới từng
  bước tinh chỉnh. Thuật toán 2-2 tinh chỉnh quá sớm.
\end{enumerate}

Thuật toán mới như sau. Dùng $a[1]$ và $a[2]$ lần lượt so sánh với
$a[3],\ldots,a[n]$, trong đó $a[i]$ biểu thị chỉ số của vật thứ $i$ trong vùng
có thể. Không mất tính tổng quát, giả sử $a[1]$ nằm bên trái $a[2]$. Khi đó có
thể chia tất cả vật thành ba loại: những vật nằm bên trái $a[1]$; những vật
nằm giữa $a[1]$ và $a[2]$; và những vật nằm bên phải $a[2]$. Dựa vào số lượng
từng loại, có thể biết vật có độ cứng trung bình là $a[1]$, là $a[2]$, hoặc
nằm trong một trong ba loại kia. Nếu nằm trong một loại, tiếp tục dùng phương
pháp tương tự để phân rã loại đó cho đến khi tìm được kết quả.

Gọi là ``tương tự'' vì có hai điểm khác. Thứ nhất, ở tầng sau không thể giả sử
$a[1]$ nằm bên trái $a[2]$; khi ấy các giá trị $a[1],a[2],\ldots,a[u]$ đã trở
thành chỉ số của các vật trong loại đang xét. Cần đưa vào một mốc từ tầng phân
loại trước, tức $a[1]$ hoặc $a[2]$ của tầng trước, và thực hiện một phép so
sánh để đảm bảo thứ tự của $a[1]$ và $a[2]$ nhất quán với trước đó. Thứ hai,
vị trí của vật trung bình không còn nhất thiết là $m+1$; vị trí cụ thể phải
được tính từ vị trí trung bình của tầng trước và kết quả phân loại. Vị trí này
được truyền như một tham số trong đệ quy.

Cuối cùng, để chống dữ liệu đặc biệt, mỗi lần chọn mốc không dùng cố định
$a[1]$ và $a[2]$, mà dùng $a[x]$ và $a[y]$, trong đó $x,y$ là các số ngẫu
nhiên.

Thuật toán này có thể giải hoàn toàn bài toán.

\section{Tổng kết}

Bảng sau so sánh hiệu suất của ba thuật toán. Mỗi kết quả là trung bình của
$10$ bộ dữ liệu đầu vào sinh ngẫu nhiên.
\[
\begin{array}{c|ccc}
& \text{2-1} & \text{2-2} & \text{3-1}\\
\hline
\text{Số lần so sánh khi } N=1449 & 11176 & 8189 & 3442
\end{array}
\]

Qua bài này, tác giả nhận ra rằng khi cải tiến một chương trình, cần hiểu đầy
đủ điểm yếu của chương trình đó và nắm đúng mâu thuẫn chính. Đồng thời, khi
thấy một thuật toán không còn nhiều không gian cải tiến, vẫn cần biết tổng kết
kinh nghiệm và những điểm thành công của nó để làm nền cho thuật toán mới.

\section*{Phụ lục}

\begin{tabular}{ll}
\texttt{medium-hardness-report.doc} & bài báo cáo này\\
\texttt{median1.pas} & chương trình của thuật toán 2-1\\
\texttt{median2.pas} & chương trình của thuật toán 2-2\\
\texttt{median3.pas} & chương trình của thuật toán 3-1\\
\texttt{medmake.pas} & chương trình sinh dữ liệu đầu vào\\
\texttt{device.pas} & đơn vị tương tác\\
\end{tabular}

\end{document}
